\documentclass[twocolumn]{aastex631}

\usepackage{amsmath}
\usepackage{multirow}
\usepackage{natbib}
\usepackage{graphicx} 
\usepackage{aas_macros}

\begin{document}

Generalized Proca theories, which extend General Relativity with vector and scalar fields, frequently suffer from ghost and tachyon instabilities that can re-emerge at the quantum level, posing a significant challenge to their theoretical viability. To address this, we introduce a novel dynamically-coupled disformal-Galileon symmetry, unifying a Galilean transformation of the scalar field with a field-dependent disformal transformation of the metric, where disformal factors are explicitly constructed from scalar Galileon-invariant terms. We define explicit transformation rules for the scalar, vector, and metric fields within the Generalized Proca Lagrangian, ensuring its invariance and providing robust theoretical stability. Through rigorous tree-level analysis on a Minkowski background, we confirm ghost and tachyon freedom for all five propagating degrees of freedom. Crucially, employing the background field method and associated Ward identities, we demonstrate this symmetry acts as a non-renormalization theorem, preventing the radiative generation of ghost-inducing operators at the one-loop level and confirming the theory's technical naturalness as an effective field theory. We further investigate the theory's compatibility with cosmological observations, showing that the stringent gravitational wave speed constraint can be satisfied by appropriate model specifications. While an initial numerical cosmological analysis for a specific parameter set, performed by integrating the modified Friedmann and field equations, did not quantitatively reproduce the observed dark energy density and equation of state, it confirmed the model's testability and highlighted the necessity for comprehensive parameter space analysis. This work establishes the dynamically-coupled disformal-Galileon symmetry as a powerful and robust mechanism for constructing quantum-stable and observationally testable theories of modified gravity.

\end{document}
                